% Compile with:
%   pdflatex reports/transchrombp_paper_outline.tex

\documentclass[11pt]{article}

\usepackage[a4paper,margin=1in]{geometry}
\usepackage[colorlinks=true,linkcolor=blue,urlcolor=blue,citecolor=blue]{hyperref}
\usepackage{xcolor}
\usepackage{enumitem}

\newcommand{\todo}[1]{\textcolor{red}{[TODO: #1]}}

\title{TransChromBP: Bridging AlphaGenome-Style Long-Context Modeling and ChromBPNet-Style Bias Factorization for Base-Resolution ATAC-seq Prediction}
\author{Author List TBD}
\date{Restructured Draft Outline --- March 15, 2026}

\begin{document}

\maketitle

\begin{abstract}
This paper should no longer be framed as ``a ChromBPNet variant with a
Transformer'' or as a report centered on data-processing tweaks. The real
question is whether three ingredients interact productively in base-resolution
ATAC-seq modeling: long-context sequence encoding, teacher-student
distillation, and ChromBPNet-style bias factorization. Our revised thesis is
that the key gap between AlphaGenome and ChromBPNet is not merely context
length, but the missing joint study of long-range representation learning,
knowledge transfer, and assay-specific bias decomposition within a single ATAC
prediction task. The paper will therefore present a flagship model that combines
a convolutional local encoder, a compressed-token Transformer, an explicit bias
branch, and distillation from a stronger teacher to a lighter student, all
under a fixed ATAC-specific data semantics. The goal is twofold: to establish a
strong task-specific model and to explain how Transformer, distillation, and
bias factorization contribute individually and jointly.
\end{abstract}

\tableofcontents

\section{Positioning}

\subsection{The Gap We Actually Care About}
The motivating gap is not simply that ChromBPNet is shorter-context and
AlphaGenome is longer-context. The real gap is that ATAC profile prediction
lacks a clean experimental framework to study:
\begin{itemize}[leftmargin=2em]
  \item whether long-range context helps beyond larger model capacity,
  \item whether distillation can preserve those gains in a lighter model, and
  \item whether explicit bias factorization remains necessary once the first two
  ingredients are added.
\end{itemize}

\subsection{Core Claim}
The paper should argue the following:
\begin{quote}
For base-resolution ATAC-seq profile prediction, the right design point is not
simply more context and not a bias-agnostic foundation-model formulation, but a
controlled combination of Transformer-based long-context modeling,
teacher-student distillation, and ChromBPNet-style bias factorization.
\end{quote}

\subsection{What The Paper Must Achieve}
This paper needs both of the following:
\begin{enumerate}[leftmargin=2em]
  \item a sufficiently strong flagship model that can plausibly define the best
  task-specific result on the benchmark we use;
  \item a mechanism study around that flagship model, rather than a list of
  disconnected ablations.
\end{enumerate}

\subsection{What We Are Not Claiming}
To keep the framing defensible:
\begin{itemize}[leftmargin=2em]
  \item we are not presenting a universal multimodal foundation model,
  \item we are not claiming to reproduce the full AlphaGenome scope,
  \item we are not reducing gains to context length alone,
  \item we are not turning this paper back into a data-semantics-only story.
\end{itemize}

\section{Research Questions}

The experiments should be organized around four explicit research questions:
\begin{enumerate}[leftmargin=2em]
  \item \textbf{RQ1:} Does the Transformer help because it captures useful
  long-range context, rather than merely increasing capacity?
  \item \textbf{RQ2:} Can distillation transfer teacher knowledge to a lighter
  student without losing fine-grained profile quality?
  \item \textbf{RQ3:} Once Transformer and distillation are present, is
  ChromBPNet-style bias factorization still necessary, and if so, should the
  bias fusion be fixed or learnable?
  \item \textbf{RQ4:} Are these three ingredients additive, redundant, or
  strongly interacting?
\end{enumerate}

\section{Fixed Task Semantics}

\subsection{Why Data Semantics Should Be Fixed, Not Reopened}
Current internal evidence already supports three conclusions:
\begin{itemize}[leftmargin=2em]
  \item the existing \texttt{bias\_pretrain -> main\_learnable} line is stronger
  than the current baseline on held-out test data;
  \item the bundled \texttt{hybrid\_data} variant behaves more like a useful
  negative control than a new production line;
  \item therefore, data semantics should now be treated as a fixed experimental
  condition rather than the main scientific variable.
\end{itemize}

\subsection{Task Semantics To Keep Fixed}
The paper should present the following as the default ATAC task definition:
\begin{itemize}[leftmargin=2em]
  \item shifted unstranded insertion bigWig as the canonical target signal,
  \item peak-centered supervised windows,
  \item GC-matched nonpeak negatives,
  \item nonpeak-only bias pretraining with no jitter,
  \item peak-centric main training with limited background exposure,
  \item train-time reverse-complement augmentation only,
  \item automatic count-loss calibration rather than a fixed constant.
\end{itemize}

This lets the paper focus on the intended method question: architecture,
distillation, and bias interaction.

\section{Method}

\subsection{Flagship Model}
The flagship model can be summarized as:
\begin{quote}
\textbf{Teacher-Student TransChromBP}: a convolutional local encoder, a
compressed-token Transformer, a ChromBPNet-style bias branch, profile/count
heads, and knowledge distillation from a stronger teacher to a lighter student.
\end{quote}

\subsection{Why A Teacher-Student Design Matters}
Without distillation, the story risks collapsing into ``a larger model is
better.'' With teacher-student training, we can ask whether the useful part of
the teacher is transferable and whether the transferred knowledge changes the
role of the bias branch.

\subsection{Modules}
The method section should describe:
\begin{itemize}[leftmargin=2em]
  \item a convolutional stem for motif-scale and short-range syntax,
  \item a Transformer for longer-range context,
  \item a bias branch for assay-specific cleavage preferences,
  \item a fusion layer that combines debiased and bias-derived predictions,
  \item profile and count heads,
  \item teacher and student variants that share the same ATAC task semantics.
\end{itemize}

\subsection{Distillation Targets}
Distillation should include at least:
\begin{itemize}[leftmargin=2em]
  \item output distillation on profile logits / distributions and log-count,
  \item optional representation distillation on pooled or token-level features.
\end{itemize}

\subsection{Bias-Aware Fusion}
This subsection must clearly distinguish:
\begin{itemize}[leftmargin=2em]
  \item the debiased regulatory branch,
  \item the bias branch,
  \item the full prediction after fusion,
  \item fixed versus learnable bias injection strength.
\end{itemize}

\section{Experimental Program}

\subsection{Benchmark}
The main benchmark should start from the official ChromBPNet tutorial data,
because it already gives:
\begin{itemize}[leftmargin=2em]
  \item a task-matched baseline,
  \item a reproduced ChromBPNet reference,
  \item a verified TransChromBP real-data backend,
  \item a held-out evaluation protocol we already trust.
\end{itemize}
\todo{Decide whether to add a second dataset in the main text or appendix.}

\subsection{Main Comparison Set}
The primary results table should be centered on the flagship model and its most
informative controls:
\begin{enumerate}[leftmargin=2em]
  \item ChromBPNet or the strongest conv-only ATAC baseline,
  \item Student TransChromBP without distillation,
  \item Student TransChromBP with distillation,
  \item Transformer + bias, no distillation,
  \item Transformer + distillation, no bias,
  \item full flagship: Transformer + distillation + learnable bias,
  \item full flagship with fixed bias.
\end{enumerate}

\subsection{Mechanism Axes}
The experiments should explicitly isolate three axes:
\begin{itemize}[leftmargin=2em]
  \item \textbf{Backbone axis}: conv-only vs conv+Transformer,
  \item \textbf{Training axis}: direct training vs distillation,
  \item \textbf{Bias axis}: no bias vs fixed bias vs learnable bias.
\end{itemize}

\subsection{Teacher Design}
The teacher does not need to be a full AlphaGenome-style system. A better
starting point is a stronger same-task expert:
\begin{itemize}[leftmargin=2em]
  \item longer input length,
  \item larger hidden size or more Transformer blocks,
  \item the same ATAC-specific targets and bias-aware formulation as the student.
\end{itemize}

\subsection{Metrics and Interpretability}
The paper should report:
\begin{itemize}[leftmargin=2em]
  \item overall, peak, and nonpeak metrics,
  \item profile and count metrics,
  \item full versus debiased outputs,
  \item locus-level visualizations,
  \item if possible, SHAP / motif / MoDISco analyses,
  \item if mature enough, ATAC-specific scorers such as center-mask or interval
  scorers.
\end{itemize}

\section{How Current Evidence Supports The Paper}

Current internal results already give useful constraints:
\begin{itemize}[leftmargin=2em]
  \item the existing bias-pretrained main line is a credible starting point,
  \item data semantics should be fixed instead of continually reopened,
  \item bias-only sanity checks support the idea that background bias is
  learnable but should not replace the regulatory task.
\end{itemize}

These findings should be used to justify the new paper direction, not to define
the final contribution.

\section{Expected Contributions}

The contribution list should end up close to the following:
\begin{enumerate}[leftmargin=2em]
  \item We propose an ATAC-specific flagship model that combines long-context
  Transformer modeling, teacher-student distillation, and explicit bias
  factorization.
  \item We establish a stronger benchmark model under fixed ATAC task semantics.
  \item We systematically characterize how Transformer, distillation, and bias
  factorization interact in base-resolution ATAC prediction.
  \item We show, through stratified and interpretability analyses, whether gains
  arise from improved regulatory modeling, improved bias handling, or their
  interaction.
\end{enumerate}

\section{Writing Discipline}

The main writing risk is losing focus again. Three rules should remain fixed:
\begin{itemize}[leftmargin=2em]
  \item this is an ATAC paper, not a generic genomics foundation-model paper;
  \item this is a strong-model-plus-mechanism paper, not a loose ablation dump;
  \item this is not a renewed debate about data processing as the central thesis.
\end{itemize}

Every experiment should therefore be judged by one criterion:
\begin{quote}
Does it strengthen either the flagship model or the explanation of how
Transformer, distillation, and bias factorization interact?
\end{quote}

\end{document}
